\section{固结在动系上的点的相对速度}

设坐标系$i$（$O\vec{e}_i$）为参考系，坐标系$j$（$O\vec{e}_j$）为动系，点$P$固结在动系$j$上。点$P$在$j$系中的坐标列阵$\boldsymbol{r}_j$为常值，其在$i$系中的坐标列阵为
\begin{equation}
\boldsymbol{r}_i = A_{ij}\boldsymbol{r}_j
\end{equation}
对时间求导，因$\boldsymbol{r}_j$为常值，$\dot{\boldsymbol{r}}_j = \boldsymbol{0}$，故
\begin{equation}
\dot{\boldsymbol{r}}_i = \dot{A}_{ij}\boldsymbol{r}_j + A_{ij}\dot{\boldsymbol{r}}_j = \dot{A}_{ij}\boldsymbol{r}_j
\end{equation}

\section{两右手笛卡尔直角坐标系的相对方位}

设坐标系$i$和坐标系$j$均为右手笛卡尔直角坐标系，二者的相对方位由方向余弦矩阵$A_{ij}$描述：
\begin{equation}
\vec{e}_i = A_{ij}\vec{e}_j
\end{equation}
方向余弦矩阵$A_{ij}$为正交矩阵，满足
\begin{equation}
A_{ij}^{-1} = A_{ij}^{\mathrm{T}}
\end{equation}
由此可得
\begin{equation}
I_3 = A_{ij}A_{ij}^{\mathrm{T}}, \quad I_3 = A_{ij}^{\mathrm{T}}A_{ij}
\label{eq:6.1}
\end{equation}

\section{方向余弦矩阵的时间变化率（由$I_3 = A_{ij}A_{ij}^{\mathrm{T}}$出发）}

对$I_3 = A_{ij}A_{ij}^{\mathrm{T}}$两端同时求导：
\begin{equation}
O_3 = \frac{\mathrm{d}}{\mathrm{d}t}\left(A_{ij}A_{ij}^{\mathrm{T}}\right) = \dot{A}_{ij}A_{ij}^{\mathrm{T}} + A_{ij}\dot{A}_{ij}^{\mathrm{T}}
= \left(\dot{A}_{ij}A_{ij}^{\mathrm{T}}\right) + \left(\dot{A}_{ij}A_{ij}^{\mathrm{T}}\right)^{\mathrm{T}}
\end{equation}
令
\begin{equation}
\dot{A}_{ij}A_{ij}^{\mathrm{T}} =: \tilde{\boldsymbol{b}}
\end{equation}
则
\begin{equation}
\dot{A}_{ij} = \tilde{\boldsymbol{b}}A_{ij}
\end{equation}
其中$\tilde{\boldsymbol{b}}$为反对称矩阵。

\section{方向余弦矩阵的时间变化率（由$I_3 = A_{ij}^{\mathrm{T}}A_{ij}$出发）}

对$I_3 = A_{ij}^{\mathrm{T}}A_{ij}$两端同时求导：
\begin{equation}
O_3 = \frac{\mathrm{d}}{\mathrm{d}t}\left(A_{ij}^{\mathrm{T}}A_{ij}\right) = \dot{A}_{ij}^{\mathrm{T}}A_{ij} + A_{ij}^{\mathrm{T}}\dot{A}_{ij}
= \left(A_{ij}^{\mathrm{T}}\dot{A}_{ij}\right)^{\mathrm{T}} + \left(A_{ij}^{\mathrm{T}}\dot{A}_{ij}\right)
\end{equation}
令
\begin{equation}
A_{ij}^{\mathrm{T}}\dot{A}_{ij} =: \tilde{\boldsymbol{d}}
\end{equation}
则
\begin{equation}
\dot{A}_{ij} = A_{ij}\tilde{\boldsymbol{d}}
\end{equation}

\section{方向余弦矩阵的时间变化率——两种推导的统一}

综合6.3节和6.4节的结果，方向余弦矩阵的时间变化率满足
\begin{equation}
\begin{cases}
\dot{A}_{ij}A_{ij}^{\mathrm{T}} = \tilde{\boldsymbol{b}}, \\[4pt]
A_{ij}^{\mathrm{T}}\dot{A}_{ij} = \tilde{\boldsymbol{d}},
\end{cases}
\quad\Rightarrow\quad
\begin{cases}
\dot{A}_{ij} = \tilde{\boldsymbol{b}}A_{ij}, \\[4pt]
\dot{A}_{ij} = A_{ij}\tilde{\boldsymbol{d}},
\end{cases}
\quad\Rightarrow\quad
\tilde{\boldsymbol{b}}A_{ij} = A_{ij}\tilde{\boldsymbol{d}}
\end{equation}

结合方向余弦矩阵的坐标变换性质可知：$\boldsymbol{b}$和$\boldsymbol{d}$应为某一个与坐标系$i$和坐标系$j$相对方位有关的矢量在$i$系和$j$系中的投影列阵。该矢量可记为
\begin{equation}
\vec{\omega}_{ij} = \vec{e}_i^{\mathrm{T}}\boldsymbol{b} = \vec{e}_j^{\mathrm{T}}\boldsymbol{d}
\end{equation}
为便于记忆，作如下变量替换：
\begin{equation}
\boldsymbol{b} =: \boldsymbol{\omega}_{ij|i}, \quad \boldsymbol{d} =: \boldsymbol{\omega}_{ij|j}
\end{equation}

\section{方向余弦矩阵的时间变化率$\dot{A}_{ij}$}

方向余弦矩阵的时间变化率满足
\begin{equation}
\begin{cases}
\dot{A}_{ij}A_{ij}^{\mathrm{T}} = \tilde{\boldsymbol{\omega}}_{ij|i}, \\[4pt]
A_{ij}^{\mathrm{T}}\dot{A}_{ij} = \tilde{\boldsymbol{\omega}}_{ij|j},
\end{cases}
\quad\Rightarrow\quad
\begin{cases}
\dot{A}_{ij} = \tilde{\boldsymbol{\omega}}_{ij|i}A_{ij}, \\[4pt]
\dot{A}_{ij} = A_{ij}\tilde{\boldsymbol{\omega}}_{ij|j},
\end{cases}
\quad\Rightarrow\quad
\begin{cases}
\tilde{\boldsymbol{\omega}}_{ij|i}A_{ij} = A_{ij}\tilde{\boldsymbol{\omega}}_{ij|j}, \\[4pt]
\vec{\omega}_{ij} = \vec{e}_i^{\mathrm{T}}\tilde{\boldsymbol{\omega}}_{ij|i} = \vec{e}_j^{\mathrm{T}}\tilde{\boldsymbol{\omega}}_{ij|j}
\end{cases}
\label{eq:6.2}
\end{equation}

固结在动系上的点$P$（在$j$系中坐标为常值）的相对速度列阵可写为：
\begin{equation}
\dot{\boldsymbol{r}}_i = \dot{A}_{ij}\boldsymbol{r}_j
= \tilde{\boldsymbol{\omega}}_{ij|i}A_{ij}\boldsymbol{r}_j
= A_{ij}\tilde{\boldsymbol{\omega}}_{ij|j}\boldsymbol{r}_j
\end{equation}

固结在动系上的点的相对速度矢量：
\begin{equation}
\vec{v} \triangleq \vec{e}_i^{\mathrm{T}}\dot{\boldsymbol{r}}_i = \vec{\omega}_{ij} \times \vec{r}
\label{eq:6.3}
\end{equation}

满足$\vec{v} \perp \vec{\omega}_{ij}$，$\vec{v} \perp \vec{r}$，且
\begin{equation}
|\vec{v}| = |\vec{\omega}_{ij}|\,|\vec{r}|\,\sin\langle\vec{\omega}_{ij},\vec{r}\rangle = |\vec{\omega}_{ij}|\,d
\end{equation}
当$\vec{r} = \alpha\vec{\omega}_{ij}$时，$\vec{v} = \vec{0}$。

由此可知$j$系相对于$i$系作瞬时定轴转动，相对转动角速度为$\vec{\omega}_{ij}$。

\section{科里奥利公式}

考虑任意矢量$\vec{r}$，其坐标列阵为
\begin{equation}
\vec{r} = \vec{e}_i^{\mathrm{T}}\boldsymbol{r}_i = \vec{e}_j^{\mathrm{T}}\boldsymbol{r}_j, \quad \boldsymbol{r}_i = A_{ij}\boldsymbol{r}_j
\end{equation}

矢量$\vec{r}$坐标列阵的时间变化率：
\begin{equation}
\dot{\boldsymbol{r}}_i = A_{ij}\dot{\boldsymbol{r}}_j + \dot{A}_{ij}\boldsymbol{r}_j
= A_{ij}\dot{\boldsymbol{r}}_j + A_{ij}\tilde{\boldsymbol{\omega}}_{ij|j}\boldsymbol{r}_j
\end{equation}

即
\begin{equation}
\vec{e}_i^{\mathrm{T}}\dot{\boldsymbol{r}}_i = \vec{e}_i^{\mathrm{T}}A_{ij}\dot{\boldsymbol{r}}_j + \vec{e}_i^{\mathrm{T}}A_{ij}\tilde{\boldsymbol{\omega}}_{ij|j}\boldsymbol{r}_j
\end{equation}

由此得到科里奥利公式：
\begin{myformula}
\begin{equation}
\left.\frac{\mathrm{d}\vec{r}}{\mathrm{d}t}\right|_{i}
= \left.\frac{\mathrm{d}\vec{r}}{\mathrm{d}t}\right|_{j}
+ \vec{\omega}_{ij} \times \vec{r}
\label{eq:6.4}
\end{equation}
\end{myformula}

\section{角速度叠加定理}

考虑三个坐标系$\vec{e}_i$、$\vec{e}_j$、$\vec{e}_k$的相对方位：
\begin{equation}
\vec{e}_i = A_{ij}\vec{e}_j = A_{ij}A_{jk}\vec{e}_k, \quad A_{ik} = \vec{e}_i\cdot\vec{e}_k^{\mathrm{T}} = A_{ij}A_{jk}
\end{equation}

求导得
\begin{equation}
\dot{A}_{ik} = \dot{A}_{ij}A_{jk} + A_{ij}\dot{A}_{jk}
\end{equation}

代入方向余弦矩阵导数表达式：
\begin{equation}
A_{ik}\tilde{\boldsymbol{\omega}}_{ik|k} = A_{ij}\tilde{\boldsymbol{\omega}}_{ij|j}A_{jk} + A_{ij}A_{jk}\tilde{\boldsymbol{\omega}}_{jk|k}
\end{equation}

在$i$系中投影：
\begin{equation}
\tilde{\boldsymbol{\omega}}_{ik|i} = \tilde{\boldsymbol{\omega}}_{ij|i} + \tilde{\boldsymbol{\omega}}_{jk|i}, \quad
\boldsymbol{\omega}_{ik|i} = \boldsymbol{\omega}_{ij|i} + \boldsymbol{\omega}_{jk|i}
\end{equation}

\begin{myformula}
角速度叠加定理：
\begin{equation}
\vec{\omega}_{ik} = \vec{\omega}_{ij} + \vec{\omega}_{jk}
\label{eq:6.5}
\end{equation}
\end{myformula}
